\section{My Background}

I believe my academic training and research background have provided a strong foundation for conducting this PhD research, particularly in data analysis, energy simulation, and mathematical modelling. During my undergraduate studies, I participated in a social practice project that involved large-scale household and commercial surveys. I collected over 120 household surveys and conducted full coverage of neighbourhood shops within 12 days. The experience enhanced my skills in field data collection, data cleaning, and statistical analysis using MATLAB and Excel, which are directly relevant to the planned multi-source data integration in this research.

During my master's studies, I took courses such as Construction and Sustainable Design and Space and Facilities for School Building, where I gained hands-on experience with building energy simulation and performance evaluation. I became familiar with the use of simulation tools for analysing thermal performance, natural ventilation, and daylighting---skills that will be extended in this project through advanced modelling platforms such as ENVI-met, UWG, and EnergyPlus. Furthermore, my master's thesis focused on mathematical modelling and simulation of Braess' Paradox, which involved nonlinear system optimisation and network equilibrium analysis. This research experience honed my ability to handle complex coupled systems and develop computational models---an essential competence for establishing the bidirectional coupling framework in this PhD project.
